\begin{abstract}
Models that predict antimicrobial resistance from bacterial genomes are usually evaluated on random
splits, where related strains fall on both sides. We built a three-view stacked model (TriLens) for
6,962 \textit{Escherichia coli} genomes and seven antibiotics, and used it, with simpler models and
a determinant rule, to dissect what such an evaluation measures. On a random split the model
reached an average precision of 0.850, ahead of every external baseline; strain metadata and
assembly statistics alone reached 0.705. Holding out whole sequence types lowered it to 0.645, and
correcting for the prevalence shift a holdout causes left a true gap of 0.118: three drugs lost at
most 0.026 and four lost 0.152--0.239, with nothing between, while very major errors rose from 24\%
to 56\%. In pre-registered experiments, replacing the pan-genome catalogue with each strain's
AMRFinderPlus calls raised normalised average precision on unseen lineages from 0.450 to 0.726 for
drugs with acquired determinants. Assembly size predicted resistance within lineages, through
accessory gene content rather than batch, and carried no signal in 10,232 external isolates, where
frozen models did transfer and a sparse logistic regression on determinant calls beat the rule by
0.104 while gradient boosting transferred worse. In \textit{Campylobacter} no model beat the rule.
Very major error rates exceeded regulatory benchmarks for every approach.
\end{abstract}
